\uuid{J8z9}
\exo7id{4462}
\titre{exo7 4462}
\auteur{quercia}
\organisation{exo7}
\datecreate{2010-03-14}
\isIndication{false}
\isCorrection{false}
\chapitre{Série numérique}
\sousChapitre{Autre}
\cpge{mpsi-mp2i.nombres-reels-suites.c, pcsi.nombres-reels-suites.b, ptsi.nombres-reels-suites.b, mpsi-mp2i.nombres-reels-suites.d, pcsi.nombres-reels-suites.c, ptsi.nombres-reels-suites.c, mpsi-mp2i.nombres-reels-suites.e, pcsi.nombres-reels-suites.d, ptsi.nombres-reels-suites.d, mpsi-mp2i.procedes-sommatoires-discrets.a, pcsi.series-numeriques.a, ptsi.series-numeriques.a, mpsi-mp2i.analyse-asymptotique.c, pcsi.analyse-asymptotique.c, ptsi.analyse-asymptotique.c, bcpst1.suites-reelles, mp-mpi.series-numeriques-vectorielles.a, pt.series-numeriques.b, bcpst2.series-reelles}
\module{Analyse}
\niveau{L1}
\difficulte{}

\contenu{
\texte{
Soit $(x_n)$ une suite définie par : $x_0 > 0$,
$\forall\ n \in \N,\ x_{n+1} = x_n + x_n^2$.
}
\begin{enumerate}
    \item \question{Montrer que $x_n \to +\infty$ lorsque $n\to\infty$.}
    \item \question{On pose $u_n = 2^{-n}\ln x_n$. Montrer que la suite $(u_n)$ est convergente.
    (On étudiera la série $\sum u_{n+1}-u_n$)}
    \item \question{En déduire qu'il existe $\alpha > 0$ tel que $x_n \sim \alpha^{2^n}$.}
\end{enumerate}
}
